\section{回帰 : 関係を直線で表し, 予測に使う}

\subsection{問い : 製作費を増やしたら, 視聴回数はどのくらい変わるか}

前章では, 製作費と視聴回数のあいだに, 弱いながらも正の相関があることを相関係数で確かめた.
製作費の大きい作品ほど, 視聴回数も多い傾向にある.

この相関係数を手がかりに, もう一歩踏み込めないだろうか.
たとえば, ある作品の製作費を1千万円積み増したとき, 視聴回数がどれだけ伸びるのかを見積もりたい.
ところが, 相関係数をどう眺めても, その数字は出てこない.
相関係数は$-1$から$1$のあいだの, 単位のない数だ.
二つがどれくらい一緒に動くかは表しても, 片方を1つ動かしたときにもう片方がいくつ動くか, という量までは含んでいない.

そこで, 前章で描いた散布図に戻ってみる.
製作費と視聴回数の点が, 右上がりにばらついている.
この点の並びを, 1本の直線でなぞってみる.
すると, その傾きが, 製作費が1つ増えるごとに視聴回数が平均していくつ増えるかにあたる.
相関係数に足りなかったのは, この傾きの情報だったわけだ.

では, その1本の直線は, どうやって引けばいいのだろう.


\subsection{素朴な方法と限界 : 目で引いた線は当てになるか}

散布図に直線を1本引くだけなら, 定規があれば足りそうだ.
ところが, 実際に定規を当ててみると, 意外なところで手が止まる.
点は少しずつずれた場所に散らばっていて, すべての点の上を通る直線は引けない.
どんな線を引いても, どこかの点は線から外れる.
右上の点に寄せれば左下の点から離れるし, 全体の真ん中を狙えば, どの点にも微妙に合わない.
どの外れを我慢して, どの点に寄せるか.
それを決めるものがないまま, なんとなくの1本を引くことになる.

このなんとなくは, 人によって違う.
同じ散布図でも, 右端に離れた点を重視する人は急な線を引き, 中央の密集に合わせる人はゆるやかな線を引く.

つまり, 目分量の線には2つの問題がある.
「良い直線」とは何かの\textbf{基準がない}こと.
そして, 基準がないせいで誰が引くかによって線が変わり, \textbf{再現性がない}ことだ.

そこで, 「良い直線」とは何かの基準を, 先に決めることにする.
基準が決まれば, どの1本が最も良いかは目分量ではなく計算で決まり, 誰がやっても同じ線になる.


\subsection{外れの大きさで直線を選ぶ : 残差と最小二乗法}

\subsubsection*{直線の式 : \texorpdfstring{$y = a + bx$}{y = a + bx}}

2変数$x$と$y$の間に直線的な関係があると想定して,
\[
  y = a + bx
\]
という形でその関係を表すことを考える.

$b$は\textbf{傾き}だ.
$x$が1増えたとき, $y$が平均的にいくつ変わるかを表す.
製作費と視聴回数の例なら, $b$ は「製作費が1単位増えたときに視聴回数が平均何単位増えるか」を表す.

$a$は\textbf{切片}で, $x = 0$のときの$y$の値にあたる.
ただし, $x = 0$が現実的な意味を持たない場合も多い.
製作費がゼロの作品が実際にはありえないのなら, 切片の数値そのものにこだわる必要はない.
直線の位置を決めるために数学的に必要な部品だと思えばよい.

問題は, この$a$と$b$をどう決めるかだ.

\subsubsection*{残差 : 外れを数にする}

前の節で決め手を欠いたのは, 点の「外れ」を目分量のまま扱っていたからだ.
そこでまず, 外れを数にする.
$i$番目のデータ点$(x_i, y_i)$に対して, 直線が予測する値は$a + bx_i$だ.
実際の値$y_i$とこの予測値の差を\textbf{残差}と呼び, $e_i$と書く.

\begin{defi}{残差}{residual}
データ点$(x_i, y_i)$と回帰直線$y = a + bx$に対し,
\[
  e_i = y_i - (a + bx_i)
\]
を$i$番目の\textbf{残差}と呼ぶ.
\end{defi}

$e_i$が正なら, 実際の値が直線より上にある.
負なら, 直線より下にある.
残差は, 直線がその点をどれだけ外したかの記録だ.
すべての点を通る直線が引けない以上, 残差をゼロにはできない.
ならば, 残差がなるべく小さい直線を「良い直線」と呼ぶことにすればよさそうだ.

では, $n$個ある残差が「全体として小さい」とはどういう状態だろうか.
残差をそのまま合計すると, 正と負が打ち消し合ってゼロに近づいてしまう.
第2章で分散を定義したときも, 偏差の合計がゼロになるのを防ぐために二乗した.
ここでも同じ手を使う.
各残差を\textbf{二乗}してから合計する.

\subsubsection*{最小二乗法 : 残差の二乗和を最小にする}

この「残差の二乗和を最小にする$a$と$b$を選ぶ」という方針に名前がついている.

\begin{defi}{最小二乗法}{ols}
$n$個のデータ$(x_1, y_1), \ldots, (x_n, y_n)$に対し,
\[
  S(a, b) = \sum_{i=1}^{n} e_i^2 = \sum_{i=1}^{n} (y_i - a - bx_i)^2
\]
を最小にする$a$と$b$を選ぶ方法を\textbf{最小二乗法}（OLS: Ordinary Least Squares）と呼ぶ.
\end{defi}

$S(a, b)$は残差の二乗和だ.
$a$と$b$の値が変われば$S$の値も変わる.
探すのは, $S$を最小にする$a$と$b$の組だ.

とはいえ, 動かせる文字が2つある.
高校で二次関数の最小値を求めたときは, 動く文字は1つだった.
2つ同時に相手にするのは難しそうなので, 片方を止めて1つずつ調べることにしよう.

\subsubsection*{切片が先に決まる}

まず, 傾き$b$を適当な値に固定して, 切片$a$だけを動かす.
$S$の中身を$a$に注目して並べ直すと,
\[
  S = \sum_{i=1}^{n} \bigl( (y_i - bx_i) - a \bigr)^2
\]
と読める.
$b$は止めてあるから, $y_i - bx_i$の部分は$n$個の決まった数だ.
つまり$S$は, 決まった数それぞれから$a$を引いて, 二乗して足したものになっている.
展開して$a$について整理すると,
\[
  S = n a^2 - 2\left(\sum_{i=1}^{n}(y_i - bx_i)\right) a + \sum_{i=1}^{n}(y_i - bx_i)^2
\]
となり, $S$は$a$の二次関数だとわかる.
しかも$a^2$の係数は$n$で, 必ず正だ.
グラフは下に凸の放物線だから, 底がただ1か所あって, そこで$S$は最小になる.
底の位置は平方完成で求められる.
二次関数$pt^2 - 2qt + r$（$p > 0$）は, 平方完成すると$p\left(t - \frac{q}{p}\right)^2 + (\mbox{定数})$の形になるから, 底は$t = \frac{q}{p}$にある.
当てはめると,
\[
  a = \frac{1}{n}\sum_{i=1}^{n}(y_i - bx_i) = \frac{1}{n}\sum_{i=1}^{n}y_i - b \cdot \frac{1}{n}\sum_{i=1}^{n}x_i = \bar{y} - b\bar{x}
\]
のとき$S$は最小になる.
最も良い切片は, $n$個の数$y_i - bx_i$の平均だったわけだ.

この式は, 見た目以上のことを言っている.
$a = \bar{y} - b\bar{x}$を直線の式$y = a + bx$に戻すと, $x = \bar{x}$のとき$y = \bar{y}$になる.
つまり, 傾きをどう選ぼうと, 最良の切片を選んだ直線は必ず点$(\bar{x}, \bar{y})$, データの重心にあたる点を通る.
最良の直線の候補は, 重心を通る直線だけに絞られた.
残る問題は, 重心をどんな傾きで通り抜けるかだけだ.

\subsubsection*{傾きの公式を導く}

$a = \bar{y} - b\bar{x}$を$S$に代入して整理すると,
\[
  S = \sum_{i=1}^{n} \bigl( (y_i - \bar{y}) - b(x_i - \bar{x}) \bigr)^2
\]
となる.
残ったのは, $y$の偏差と$x$の偏差, そして傾き$b$だけの式だ.
各項の二乗を開いて, $b$の次数ごとにまとめる.
二乗を開いて部品を調べるこの手は, 前章で相関係数が$\pm 1$に収まる理由を確かめたときにも使った.
\[
  S = \left( \sum_{i=1}^{n}(x_i - \bar{x})^2 \right) b^2
      - 2\left( \sum_{i=1}^{n}(x_i - \bar{x})(y_i - \bar{y}) \right) b
      + \sum_{i=1}^{n}(y_i - \bar{y})^2
\]
ごつい見た目になったが, 3つの係数はどれもデータから計算できる決まった数で, 動くのは$b$だけだ.
$S$はまたしても二次関数になった.
$b^2$の係数は偏差の二乗の和だから正で, グラフは下に凸の放物線だ.
さきほどの$pt^2 - 2qt + r$と同じ形で, $p$にあたるのが$\sum(x_i - \bar{x})^2$, $q$にあたるのが$\sum(x_i - \bar{x})(y_i - \bar{y})$になっている.
底は$t = q/p$, つまり
\[
  b = \frac{\displaystyle\sum_{i=1}^{n}(x_i - \bar{x})(y_i - \bar{y})}{\displaystyle\sum_{i=1}^{n}(x_i - \bar{x})^2}
\]
のとき$S$は最小になる.

まとめると, 残差の二乗和を最小にする直線はただ1本で, その傾きと切片は
\[
  b = \frac{\displaystyle\sum_{i=1}^{n}(x_i - \bar{x})(y_i - \bar{y})}{\displaystyle\sum_{i=1}^{n}(x_i - \bar{x})^2}, \qquad a = \bar{y} - b\bar{x}
\]
と, データから直接計算できる.

傾き$b$の分数の部品は, 前章で出会ったものばかりだ.
分子の$\sum(x_i - \bar{x})(y_i - \bar{y})$は, $x$と$y$の偏差の積の合計だ.
前章の共分散の分子と同じ構造をしている.
分母の$\sum(x_i - \bar{x})^2$は$x$の偏差の二乗和で, $x$の分散の分子にあたる.

つまり傾き$b$は, 「$x$と$y$がどれだけ一緒に動くか」を「$x$自身の散らばり」で割ったものだ.
$x$が1単位散らばるごとに, $y$が平均的にどれだけ動くか.
公式はそう読める.

ここで, 前の節で挙げた2つの問題を思い出そう.
「良い直線」の基準がないことと, 誰が引くかで線が変わってしまうことだった.
「残差の二乗和を最小にする」という基準は, 人間が選んだ約束だ.
偏差の二乗で散らばりを測った第2章の分散と同じで, 自然にそう決まるわけではない.
残差を絶対値で測る方法も考えられるが, それだと$S$は二次関数にならず, いま使った平方完成の手が通らない.
二乗なら, 見てきたとおり答えが1つの式にまとまる.
そして, 基準さえ約束してしまえば, どの直線が最良かは二次関数の底として一意に決まり, そこに人間の裁量はもう入らない.
基準は人間が決め, その基準の下での答えは数学が決める.
だから, 誰が計算しても同じ1本の線が出る.

\begin{mybox}{ガウスとケレスの軌道}
1801年, イタリアの天文学者ピアッツィが小惑星ケレスを発見した.
しかし観測できたのはわずか41日間で, ケレスは太陽の向こう側に隠れてしまった.
再び姿を現したときにどこにいるかを, 不完全な観測データから予測しなければならない.
ガウスは観測値と予測値のずれの二乗和を最小にする方法で軌道を推定し, その予測は的中した.
ケレスは予測された位置で再発見された.

最小二乗法は, 「手元にあるのは不完全なデータだが, そこから最も確からしい答えを引き出したい」という切実な動機のもとで実用化された.
なお, フランスのルジャンドルが1805年に最小二乗法を先に公表しており, 優先権を巡る論争がある.
ガウスは「自分は1795年から使っていた」と主張したが, 出版が後だったため決着はついていない.
\end{mybox}


\subsubsection*{相関係数との関係}

傾き$b$の式と前章の相関係数$r$の式を見比べると, 次の関係が導かれる.
\[
  b = r \cdot \frac{s_y}{s_x}
\]

$r$は関係の強さと方向を$-1$から$1$で表す, 単位のない指標だった.
$s_y / s_x$は, $y$の散らばりが$x$の散らばりの何倍かを表している.
単位のない「連動の度合い」に散らばりの比を掛けると, 「$x$が1単位増えたとき$y$がいくつ変わるか」という単位つきの量に変わる.
前章の相関係数と本章の傾きは, 同じ直線関係を別の角度から見た量だったわけだ.


\subsection{引いた直線を疑う : \texorpdfstring{$R^2$}{R²}と残差プロット}

最小二乗法は, どんな点の散らばりに対しても, 残差の二乗和が最小になる1本を必ず選び出す.
点が直線とは呼べない並び方をしていても, 計算は止まらず, 傾きと切片の数字が出てくる.
選ばれた1本が候補の中で最も外れの小さい線であることは, さきほどの導出が保証している.
しかし, その「最も小さい外れ」が小さいとはかぎらない.
最もましな線と, 良い線は別物だ.
だから, 引いた直線がどのくらい信用できるかは, 引いたあとに確かめる必要がある.

\subsubsection*{\texorpdfstring{$R^2$}{R²} : 説明できた割合}

手がかりは, 最小化のあとに残ったものにある.
残差の二乗和$S$は, 最小にはなったが, ゼロになったわけではない.
残った$S$は, 直線がどうしても拾いきれなかった散らばりの量だ.
この残りが小さいのか大きいのかを測りたい.

ただ, $S$の値をそのまま見ても, 判断はつかない.
視聴回数を「回」で数えるか「万回」で数えるかで, 値は何桁も変わってしまうからだ.
共分散をそのままでは比べられず, 標準偏差で割って相関係数にしたのと同じ事情で, 比べる相手が要る.

そこで, 直線を引く前の外し方と比べることにする.
もし$x$の情報がまったくなかったとしたら, $y$の値をどう予測するだろうか.
使えるのは$\bar{y}$, つまり$y$の平均しかない.
どのデータ点に対しても$\bar{y}$と答えるしかないから, 予測は当然はずれる.
各点と$\bar{y}$のずれの二乗和$\sum(y_i - \bar{y})^2$が, $y$の全体の散らばりだ.
この量を$\mathrm{SS_{total}}$と呼ぶ.

回帰直線を引くと, 予測値が$\bar{y}$から$\hat{y}_i = a + bx_i$に変わる.
$x$の情報を使ったぶんだけ, 予測は改善されるはずだ.
それでも残るずれの二乗和$\sum(y_i - \hat{y}_i)^2$を$\mathrm{SS_{res}}$と呼ぶ.
中身は, 最小二乗法で最小にしたあの$S$の, 最小値そのものだ.
$\mathrm{SS_{total}}$は$x$を知らないときの誤差で, $\mathrm{SS_{res}}$は$x$を使ったあとの誤差だ.
この2つを比べて, $x$の情報で誤差がどれだけ減ったかを測る.

\begin{defi}{決定係数}{r-squared}
$y$の全変動を$\mathrm{SS_{total}} = \sum(y_i - \bar{y})^2$, 残差による変動を$\mathrm{SS_{res}} = \sum(y_i - \hat{y}_i)^2$とする.
決定係数は
\[
  R^2 = 1 - \frac{\mathrm{SS_{res}}}{\mathrm{SS_{total}}}
\]
で定義される. $\hat{y}_i = a + bx_i$は回帰直線による予測値である.
\end{defi}

$\mathrm{SS_{res}} / \mathrm{SS_{total}}$は, 直線で説明しきれなかった散らばりの割合だ.
1から引けば, 「説明できた割合」に変わる.

$R^2 = 0.76$なら, $y$の散らばりの76\%が$x$との直線関係で説明でき, 残り24\%は直線では捉えきれていない.
$R^2 = 1$なら全点が直線上に乗っており, $R^2 = 0$なら$x$と$y$に線形関係がない.

ここまで読んで, 「0.76は良いのか悪いのか」と思ったかもしれない.
正直に言えば, 分野による.
物理の実験データなら$R^2 = 0.76$は低すぎるかもしれないし, 人間の行動データや社会調査なら十分に高い.
「この値以上なら合格」という普遍的な基準はない.
$R^2$は他のモデルや他のデータとの比較に使う相対的な指標だと思っておくのがよい.

単回帰では, $R^2$は相関係数$r$の二乗に一致する.
$r = 0.87$なら$R^2 = 0.87^2 \approx 0.76$だ.
前章で「相関が強い」と評価したことと, 本章で「76\%説明できた」と評価することが, 数値としてつながっている.

\subsubsection*{\texorpdfstring{$R^2$}{R²}だけでは安心できない}

$R^2$が高ければ直線がデータをよく捉えている, と言いたくなる.
しかし, $R^2$は残差の大きさを1つの数値に圧縮した要約だ.
圧縮する過程で, ずれ方のパターンが消える.
数の要約だけでは分布の形が見えないことは第2章で確かめ, 失われた形は図で取り戻すのだと第3章でやった.
同じことが, 直線への当てはまりにも起きる.

アンスコムのカルテットという有名な例がある.
4組のデータ（各11点）のすべてで, $r \approx 0.82$, $R^2 \approx 0.67$, 傾きも切片もほぼ同じだ.
数値だけ見たら, 4組とも「まあまあの直線関係」で片づけてしまうだろう.

ところが散布図を描くと, 実態はまるで違う.

% [図挿入: アンスコムのカルテットの4枚の散布図]

1組目は素直な直線関係で, 回帰直線がデータをよく捉えている.
2組目は完全な放物線で, 直線を当てはめること自体が不適切だ.
3組目はほぼ完璧な直線に外れ値が1点だけあり, その1点が傾きを歪めている.
4組目は$x$が1点だけ離れた位置にあり, その1点が見かけの相関を作り出している.

1973年にフランシス・アンスコムがこの4組を発表した意図は明確だ.
統計量を計算する前にグラフを描け, ということだ.
$R^2$や傾きの数値だけでは, データの本当の姿はわからない.

\subsubsection*{残差プロット : 残差の形を目で見る}

では, 何の図を描けばよいか.
疑っている相手は「直線の外し方」だから, 見るべきは残差そのものだ.
残差$e_i = y_i - \hat{y}_i$を, 予測値$\hat{y}_i$に対してプロットした図を\textbf{残差プロット}と呼ぶ.

なぜ横軸が$x_i$ではなく$\hat{y}_i$なのかと思うかもしれない.
単回帰では$\hat{y}_i = a + bx_i$だから, 横軸を$x_i$にしても$\hat{y}_i$にしても図の形は変わらない.
ただ, 説明変数が複数ある重回帰では$x$が1つに決まらないため, $\hat{y}$を使うのが慣例だ.
本書では単回帰だけを扱うが, 慣例に合わせておく.

直線がデータの関係をうまく捉えていれば, 残差はゼロのまわりにランダムに散らばる.
逆に, 残差に系統的なパターンが見えたら, 直線では捉えきれない構造がデータに残っている.

残差プロットで確認するのは2点だ.

\begin{itemize}
  \item \textbf{パターンがあるか}. U字型が見えたら, 本当は曲線的な関係なのに直線で無理に近似した可能性がある. 扇形に広がっていれば, 予測値が大きい領域で予測が不安定になっている.
  \item \textbf{飛び離れた点があるか}. 他の残差から大きく外れた点があれば, そのデータ点を個別に確認する.
\end{itemize}

U字型が見えた場合にどうするか.
典型的な対処は, 二次関数などの曲線モデルを検討するか, データの範囲を絞って直線が妥当な区間だけを使うことだ.
本書では直線モデルの範囲にとどめるが, 直線でうまくいかないときにも出口はある.

残差の量は$R^2$で測り, 残差の形は残差プロットで見る.
この2つが, 引いた直線を疑うときの基本の確かめになる.


\subsection{実践 : 製作費と視聴回数で問いに答える}

章の冒頭で, 製作費を積み増したら視聴回数はどれだけ伸びるのか, という問いを立てた.
前章で扱った, 配信サービスの邦画作品の標本から製作費と視聴回数の組を10件取り出して, 直線を引き, その直線を疑い, 問いに数字で答えるところまで通してみよう.
以下の表の数値は, 計算手順を見せるための仮置きであって, ランニングデータの実値とは別物である（実値はランニングデータ確定後に置き換える）.

\subsubsection*{Step 1 : 散布図の確認}

まず, データの全体像を確認する.

\begin{center}
\begin{tabular}{crr}
\hline
$i$ & 製作費 $x_i$（仮単位） & 視聴回数 $y_i$（仮単位） \\
\hline
1 & 20 & 140 \\
2 & 25 & 200 \\
3 & 30 & 215 \\
4 & 40 & 180 \\
5 & 45 & 245 \\
6 & 50 & 200 \\
7 & 55 & 260 \\
8 & 65 & 255 \\
9 & 70 & 310 \\
10 & 80 & 295 \\
\hline
\end{tabular}
\end{center}

% [図挿入: 製作費 vs 視聴回数の散布図]

散布図を描くと, 右上がりの傾向が見える.
ただし, 点はきれいに並んでいるわけではない.
たとえば$i = 4$（製作費40, 視聴回数180）は周囲より低く, $i = 5$（製作費45, 視聴回数245）は高い.
同じような製作費でも視聴回数には幅がある. 相関が1ではないというのは, 散布図の上ではこの幅のことだ.
大きな曲がりや極端に飛び離れた点は見当たらないので, 直線を当てはめることに無理はなさそうだ.

\subsubsection*{Step 2 : 傾きと切片の計算}

公式に代入するために, まず平均を求める.
\[
  \bar{x} = \frac{20 + 25 + \cdots + 80}{10} = \frac{480}{10} = 48, \qquad
  \bar{y} = \frac{140 + 200 + \cdots + 295}{10} = \frac{2300}{10} = 230
\]

次に, 各データの偏差と, その積を計算する.
面倒に見えるが, 表にすれば機械的に進む.

\begin{center}
\begin{tabular}{crrrrr}
\hline
$i$ & $x_i$ & $y_i$ & $x_i - \bar{x}$ & $y_i - \bar{y}$ & $(x_i - \bar{x})(y_i - \bar{y})$ \\
\hline
1 & 20 & 140 & $-28$ & $-90$ & 2520 \\
2 & 25 & 200 & $-23$ & $-30$ & 690 \\
3 & 30 & 215 & $-18$ & $-15$ & 270 \\
4 & 40 & 180 & $-8$ & $-50$ & 400 \\
5 & 45 & 245 & $-3$ & 15 & $-45$ \\
6 & 50 & 200 & 2 & $-30$ & $-60$ \\
7 & 55 & 260 & 7 & 30 & 210 \\
8 & 65 & 255 & 17 & 25 & 425 \\
9 & 70 & 310 & 22 & 80 & 1760 \\
10 & 80 & 295 & 32 & 65 & 2080 \\
\hline
 & & & & 計 & 8250 \\
\hline
\end{tabular}
\end{center}

同様に$\sum(x_i - \bar{x})^2$を求めると, $784 + 529 + 324 + 64 + 9 + 4 + 49 + 289 + 484 + 1024 = 3560$となる.

これで部品が揃った. 公式に代入する.
\[
  b = \frac{8250}{3560} \approx 2.32, \qquad a = 230 - 2.32 \times 48 \approx 118.6
\]

回帰直線は$y \approx 119 + 2.3x$と表せる.

% [図挿入: 散布図に回帰直線を重ねたもの]

傾き$b \approx 2.3$は, 「製作費を1単位増やすと, 視聴回数が平均的に約2.3単位増える」と読める.
相関係数だけではわからなかった「いくつ変わるか」が, この1つの数字に入っている.
冒頭の問いへの答えは, ひとまずこの形で出た.

切片$a \approx 119$は, 式の上では「製作費がゼロのとき視聴回数は約119」となる.
しかし, データの製作費は20から80の範囲に収まっている.
製作費ゼロの作品は実データに存在しない.
この119は直線の位置を決めるために計算上現れた数値であって, 「製作費を一切かけなくても119の視聴回数が出る」と解釈するのは危うい.

\subsubsection*{Step 3 : \texorpdfstring{$R^2$}{R²}で当てはまりを確認する}

最小二乗法は, どんなデータからでも1本の線を出すのだった.
出てきた数字をどこまで信じてよいかを, まず$R^2$で確かめる.

仮にこの仮置きデータで$r = 0.87$と求めたとすれば, $R^2 = 0.87^2 \approx 0.76$だ.
視聴回数の散らばりの約76\%が, 製作費との直線関係で説明できている.
残りの24\%は製作費だけでは説明しきれない.
おすすめ表示回数, 評価件数, ジャンル, 公開時期など, 他の要因が混ざっているのだろう.
表の$i = 4$（製作費40なのに視聴回数180）や$i = 6$（製作費50で視聴回数200）のように, 直線から大きくずれた作品がそれにあたる.
なお, 実際のランニングデータでは製作費と視聴回数の相関はもっと弱め（おおむね $r$ で 0.3〜0.5 程度）になる想定で, $R^2$ もこれより小さい値になる.

\subsubsection*{Step 4 : 残差プロットを確認する}

各データ点の予測値$\hat{y}_i = 119 + 2.3 x_i$と残差$e_i = y_i - \hat{y}_i$を計算すると, 次のようになる.

\begin{center}
\begin{tabular}{crrrrr}
\hline
$i$ & $x_i$ & $y_i$ & $\hat{y}_i$ & $e_i$ \\
\hline
1 & 20 & 140 & 165 & $-25$ \\
2 & 25 & 200 & 177 & 23 \\
3 & 30 & 215 & 188 & 27 \\
4 & 40 & 180 & 211 & $-31$ \\
5 & 45 & 245 & 223 & 22 \\
6 & 50 & 200 & 234 & $-34$ \\
7 & 55 & 260 & 246 & 14 \\
8 & 65 & 255 & 269 & $-14$ \\
9 & 70 & 310 & 280 & 30 \\
10 & 80 & 295 & 303 & $-8$ \\
\hline
\end{tabular}
\end{center}

% [図挿入: 残差プロット]

残差の符号は$-, +, +, -, +, -, +, -, +, -$で, 正と負が入り混じっている.
U字型のパターンや, 右に行くほど広がる扇形も見えない.
ゼロのまわりにランダムに散らばっていると判断してよい.
この直線モデルは, おおむね妥当だ.

\subsubsection*{予測 : できることと, やってはいけないこと}

疑いをひととおり通過した直線は, まだ見ていない作品の視聴回数の見積もりに使える.
「製作費が50の作品の視聴回数はおよそ$119 + 2.3 \times 50 \approx 234$」と見積もれる.
50はデータの範囲（20から80）に収まっているから, この予測にはそれなりの根拠がある.

同じ式に製作費200を入れても, 計算は止まらず, $119 + 2.3 \times 200 = 579$という数字が出てくる.
しかし, 200はデータの範囲をはるかに超えた値だ.
直線関係がそこまで続いている保証はない.
現実を考えても, 製作費には効果の頭打ちがある.
予算を2倍にしたからといって視聴回数が2倍になるとは限らない.

データの範囲外に直線を延ばして予測すること, いわゆる外挿は, 「直線が無限に続く」という仮定を暗に置いている.
その仮定が正しい保証はどこにもない.

前章の因果の議論も, ここでそのまま当てはまる.
$b \approx 2.3$は「製作費が1単位多い作品は視聴回数も平均約2.3単位多い」という関係を記述しているだけであり, 「製作費を増やせば視聴回数が増える」という因果を証明しているわけではない.
回帰にできるのは関係の記述と, その記述に基づく予測までで, 因果の証明は守備範囲の外にある.


\subsection{直線を引いた先に残る問い}

本章では, 相関係数の先にあった問い, 「$x$が変わったとき$y$はどのくらい変わるか」に答えた.

最小二乗法で直線を引き, 傾き$b \approx 2.3$から「製作費が1単位多い作品は視聴回数も平均約2.3単位多い」という量的な関係を読み取った.
$R^2 \approx 0.76$でその直線の説明力を確認し, 残差プロットで直線を信用してよいかを目で判断した.
前章では「関係がある」までだった答えが, 「1単位あたり約2.3」という数字まで進んだ.

振り返ると, 章の半分は線の引き方ではなく, 引いた線の疑い方に使った.
最小二乗法はどんなデータからでも1本の線を出すから, 出てきた数字を信じてよいかは, 残差の量と形を確かめてはじめて言える.
線を引いて終わりにせず, 残差を読むところまでが回帰の手続きだ.

ただし, 最小二乗法にも$R^2$にも前提がある.
$x$と$y$の関係がおおむね直線的であること.
各データ点が互いに独立であること.
極端な外れ値がないこと.
前提が崩れていれば, 直線を引いて数字は出るが, その数字は当てにならない.

たとえば, もしデータ点が独立でないのに独立だと仮定して回帰すると, 残差にパターンが現れ, $R^2$が実際の説明力より高く出るかもしれない.
回帰直線の数字を眺めても, そのことには気づけない.
これは直線を引く技術の問題ではなく, データの設計と解釈の問題だ.

前提が崩れる場面は意外と多い.
標本の集め方に偏りがあれば, 回帰直線も偏る.
因果の方向を取り違えていれば, 回帰の解釈そのものが間違う.
次章では, 第2章から本章まで積み上げてきた手法の前提を整理し, その前提が崩れるとどうなるかを体系的に確認する.
あわせて, バイアスの類型と, データを扱うときの倫理も扱う.
手法の「外側」を知ることが, 手法を正しく使うための最後の一歩になる.
